% ============================================================
%  00_acronymes_manuel.tex  –  Aide-mémoire usage des acronymes
% ============================================================
% Ce fichier n'est PAS inclus dans la compilation.
% Il sert de référence pour l'utilisation des macros \ac{}, \acf{}, \acp{}.
%
% Syntaxe courante :
%   \ac{JWT}   → première occurrence : "JSON Web Token (JWT)", suivantes : "JWT"
%   \acf{JWT}  → forme complète forcée : "JSON Web Token (JWT)"
%   \acs{JWT}  → forme courte forcée   : "JWT"
%   \acl{JWT}  → forme longue forcée   : "JSON Web Token"
%
% Exemples dans le texte :
%   "L'authentification repose sur \ac{JWT}."
%   "Le backend expose une \ac{API} REST."
%   "La plateforme utilise \ac{BDD} PostgreSQL."
